\chapter{Parametric Quadratic Programming}
\label{ch:parametric}

So far a QP has been a single problem with a single answer. This chapter changes
the question. Attach a scalar parameter to the linear term and ask how the solution
moves as the parameter varies:
\begin{equation}
\label{eq:pqp}
\begin{aligned}
x(\lambda) = \argmin_{x \in \R^n}\ &\tfrac12\, x\trans H x - \lambda\, d\trans x \\
\text{s.t.}\quad &A x = b,\ \ G x \le h,\ \ \ell \le x \le u .
\end{aligned}
\end{equation}
We have switched to the symbols the parametric literature uses---$H$ for the
quadratic form, $A$ and $G$ for equality and inequality systems, $\ell,u$ for
bounds---because Part~III is where this book meets four other literatures and
Appendix~\ref{app:notation} has to translate somewhere.

The answer is remarkable and is the foundation of everything in Part~III:
\emph{the solution map is piecewise linear}. On any interval of $\lambda$ where the
active set is constant, $x(\lambda)$ is an affine function of $\lambda$; the active
set changes at finitely many \emph{breakpoints}; and the entire family can therefore
be computed \emph{exactly}, corner to corner, by one finite algorithm. Not
approximated on a grid. Computed.

\section{The standing assumption}

We require only that $H \succeq 0$ and that the \emph{reduced} Hessian
$H_{\Free\Free}$ on the free block is positive definite at every active set the path
visits. This is weaker than the $H \succ 0$ used in Parts~I and~II, and deliberately
so: it is exactly what the regression instance of Chapter~\ref{ch:homotopy} needs.
There $H = X\T X$ is singular whenever the design has more columns than rows---the
very regime the LASSO exists for---yet the free block $X_\Free\T X_\Free$ is positive
definite as soon as the active columns are linearly independent, which holds in
general position. The finance instance has $H = \Sig \succ 0$ outright, so the
condition is automatic there.

\begin{remark}[Where the assumption bites]
\label{rem:rank}
This is an assumption and not a formality, and it is worth saying where it fails
rather than burying it. When $H_{\Free\Free}$ is singular at some visited active set,
the bordered system~\eqref{eq:border} below no longer determines the segment, and
what the path even \emph{is} becomes a question of which minimiser one selects.
Existing treatments divide on the answer. Tibshirani and Taylor \cite{tibshirani2011}
for the generalized lasso, and Gaines, Kim and Zhou \cite{gaines2018} for the
constrained LASSO, both restore strong convexity by adding a small ridge, so the
algorithms they analyse are those of a perturbed problem. Tibshirani
\cite{tibshirani2013}, by contrast, gives a LASSO path algorithm valid for
\emph{arbitrary} designs with no general-position hypothesis: it succeeds by tracing
the minimum-$\ell_2$-norm solution specifically, and turns on an
insertion--deletion property of that particular selection rather than on the segment
structure alone. That is a genuinely different argument.

No counterpart is known for the generalized or constrained problem. Whether the
bordered system admits a canonical selection playing the same role is, as far as we
know, open, and Chapter~\ref{ch:reading} returns to it as one of the subject's
genuine unsolved questions.
\end{remark}

\section{Affine segments}

\begin{lemma}[Affine segments and the bordered solve]
\label{lem:affine}
\index{affine segment}
On any maximal interval of $\lambda$ on which the active set is constant, the
optimiser of \eqref{eq:pqp} is affine,
\begin{equation}
\label{eq:affine}
x(\lambda) = \alpha + \lambda\,\beta ,
\end{equation}
and $(\alpha,\beta)$ come from a single bordered linear solve~\eqref{eq:border}
whose matrix is constant on the interval. An active inequality row enters that solve
identically to an equality row.
\end{lemma}

\begin{proof}
Fix $\lambda$ in the interior of the interval and name the active set: the
equalities $A$ (always active); the active inequality rows
$\mathcal{S} \subseteq \{1,\dots,p\}$, those held at $g_i\T x = h_i$, the remainder
carrying a zero multiplier by complementarity; and the blocked variables $\Bound$
pinned at a bound, with free set $\Free = \Bound^c$. Stack the active linear
constraints as $C = [A; G_{\mathcal{S}}]$ with right-hand side
$r = (b; h_{\mathcal{S}})$.

Because \eqref{eq:pqp} is convex with linear constraints, the KKT conditions are
necessary and sufficient. Stationarity is $Hx - \lambda d + C\T\zeta + \rho = 0$,
with $\zeta$ the equality and active-inequality multipliers and $\rho$ the bound
multipliers, supported on $\Bound$. Read the $\Free$ rows, where $\rho_\Free = 0$,
and split $(Hx)_\Free = H_{\Free\Free}x_\Free + H_{\Free\Bound}x_\Bound$ with
$x_\Bound$ fixed at its bound values. Primal feasibility of the active rows is
$C_\Free x_\Free + C_\Bound x_\Bound = r$. Collecting the unknowns
$(x_\Free,\zeta)$,
\begin{equation}
\label{eq:border}
\begin{bmatrix} H_{\Free\Free} & C_\Free\T \\[2pt] C_\Free & 0 \end{bmatrix}
\begin{bmatrix} x_\Free \\[2pt] \zeta \end{bmatrix}
=
\begin{bmatrix} \lambda\, d_\Free - H_{\Free\Bound}\,x_\Bound \\[2pt]
                r - C_\Bound\,x_\Bound \end{bmatrix}.
\end{equation}
The coefficient matrix is nonsingular by exactly the argument of
Proposition~\ref{prop:saddlenonsing}, using $H_{\Free\Free}\succ0$ and full row rank
of $C_\Free$. It depends only on $H$ and the active normals, hence is constant on
the interval, while the right-hand side is affine in $\lambda$---the only $\lambda$
term being $\lambda d_\Free$. So $(x_\Free(\lambda),\zeta(\lambda))$ is affine, and
appending the $\lambda$-independent blocked entries gives \eqref{eq:affine} with
$\beta_\Bound = 0$.

Finally, $G_{\mathcal{S}}$ enters~\eqref{eq:border} only as extra rows of $C$,
indistinguishable from the equality rows, and the box bounds are carried by the
partition $(\Free,\Bound)$ rather than folded into $C$. A polyhedral penalty
$\lambda\norm{x}_1$ is the same statement with $d$ replaced, on each segment, by the
constant subgradient sign pattern of the support---the only change needed to
read~\eqref{eq:border} as the LASSO segment solve.
\end{proof}

Each homotopy step is therefore \emph{one solve} of~\eqref{eq:border}: a
factorisation of $H_{\Free\Free}$ bordered by the active normals through its Schur
complement, which is exactly the block elimination~\eqref{eq:blockelim} of
Chapter~\ref{ch:freeblock}. In practice one solves it for two right-hand sides at
once against a single factorisation, obtaining $\alpha$ and $\beta$ together.

\section{How the active set evolves}
\label{sec:events}

A representation is not an algorithm. Lemma~\ref{lem:affine} says what a segment
\emph{is}; it does not say how one segment gives way to the next. That second half
is where each field's paper does its technical work, and where the instances of
Chapter~\ref{ch:homotopy} are least recognisable to one another. We write it once.

Fix the current breakpoint $\lambda_k$, the current active set, and a direction of
travel $\sigma \in \{+1,-1\}$; write $\lambda = \lambda_k + \sigma t$ and call
$t \ge 0$ the step. Solving \eqref{eq:border} twice against one factorisation gives
the segment $x_\Free(\lambda) = \alpha_\Free + \lambda\beta_\Free$ and the
multipliers $\zeta(\lambda) = \zeta^0 + \lambda\zeta^1$; the bound multipliers follow
from stationarity as $\rho(\lambda) = \lambda d - Hx(\lambda) - C\T\zeta(\lambda)$,
affine in $\lambda$ like everything else. Write $\dot\phi = \sigma\phi_1$ for the
rate of change of any affine $\phi(\lambda) = \phi_0 + \lambda\phi_1$ along the
direction of travel.

\emph{Four things, and only four, can end a segment}: a primal quantity can reach a
bound it must respect, or a dual quantity can reach the edge of the cone it must lie
in.
\index{events!the four families}

\paragraph{$\mathrm{P}_1$: a free variable reaches a bound.} For $i \in \Free$ the
coordinate must stay in $[\ell_i,u_i]$, so
\[
t^{\mathrm{P}_1}_i =
\begin{cases}
(u_i - x_i)/\dot x_i, & \dot x_i>0,\\
(\ell_i - x_i)/\dot x_i, & \dot x_i<0,\\
+\infty, & \dot x_i=0,
\end{cases}
\]
and the event moves $i$ from $\Free$ to $\Bound$, pinned at the bound it met.

\paragraph{$\mathrm{P}_2$: an inactive inequality becomes active.} For
$j \notin \mathcal{S}$ the slack $h_j - g_j\T x$ must stay non-negative, giving
$t^{\mathrm{P}_2}_j = (h_j - g_j\T x)/(g_j\T\dot x)$ when $g_j\T\dot x>0$ and
$+\infty$ otherwise; the event adjoins $j$ to $\mathcal{S}$, adding a row to $C$.

\paragraph{$\mathrm{D}_1$: a blocked variable releases.} For $i \in \Bound$ dual
feasibility confines $\rho_i(\lambda)$ to an interval
$[\underline\rho_i(\lambda), \overline\rho_i(\lambda)]$ whose endpoints are
themselves affine in $\lambda$. Both margins $e^-_i = \rho_i - \underline\rho_i$ and
$e^+_i = \overline\rho_i - \rho_i$ are affine and non-negative at $\lambda_k$, and
the step is the first zero of either,
\[
t^{\mathrm{D}_1}_i = \min\bigl\{ -e^-_i/\dot e^-_i \ \text{if}\ \dot e^-_i<0,\ \
-e^+_i/\dot e^+_i \ \text{if}\ \dot e^+_i<0 \bigr\},
\]
and the event moves $i$ from $\Bound$ to $\Free$.

This is the family worth pausing on. For a genuine box the interval is
$[0,\infty)$ at an upper bound and $(-\infty,0]$ at a lower one, and the event is
simply $\rho_i = 0$. For an $\ell_1$ penalty it is $[-\lambda,\lambda]$, and the
event is $|\rho_i| = \lambda$: an affine function meeting a \emph{moving} threshold
rather than a fixed one. \emph{That single generalisation---a feasibility boundary
that travels with $\lambda$---is the whole of what separates an ``induced'' active
set from an ``imposed'' one at the level of the algorithm.}

\paragraph{$\mathrm{D}_2$: an active inequality releases.} For $j \in \mathcal{S}$
the multiplier must satisfy $\zeta_j \ge 0$, so
$t^{\mathrm{D}_2}_j = -\zeta_j/\dot\zeta_j$ when $\dot\zeta_j<0$ and $+\infty$
otherwise; the event deletes $j$ from $\mathcal{S}$. Equality rows never generate
this event, their multipliers being unsigned.

\medskip
The step is $t^\star = \min t$ over the four families, the event attaining it is
applied, and the walk continues. This is the simplex method's ratio test, and it
inherits the simplex method's degeneracy hazard---ties at $t^\star = 0$ and
cycling---together with its remedy: a lowest-index (Bland) tie-break, which makes
the walk deterministic and finite. Because exactly one index changes status per
step, the factorisation of the bordered matrix is \emph{updated} rather than
rebuilt, by Section~\ref{sec:updating} of Chapter~\ref{ch:freeblock}.

\begin{table}[ht]
\centering
\footnotesize
\caption{The four events as two literatures name them. The algorithms differ in
which families are populated, not in the families.}
\label{tab:events}
\begin{tabular}{@{}lll@{}}
\toprule
Event & LASSO\,/\,LARS & Critical Line Algorithm \\
\midrule
$\mathrm{P}_1$ & a coefficient crosses zero (the ``leave'' move)
               & an asset reaches a bound and leaves the free set \\
$\mathrm{P}_2$ & a linear inequality on $\beta$ becomes tight
               & a group or turnover constraint becomes tight \\
$\mathrm{D}_1$ & an inactive coordinate's correlation reaches $\lambda$ (``enter'')
               & an asset's reduced cost changes sign and it re-enters \\
$\mathrm{D}_2$ & an active inequality's multiplier hits zero
               & a constraint stops binding at a turning point \\
\bottomrule
\end{tabular}
\end{table}

Table~\ref{tab:events} reads the catalogue back into two dialects. Forward least
angle regression has $\mathrm{D}_1$ only, which is exactly what makes its active set
monotone; the LASSO modification restores $\mathrm{P}_1$ and with it the leave move;
the Critical Line Algorithm runs $\mathrm{P}_1$ and $\mathrm{D}_1$ against genuine
box bounds, with fixed rather than moving feasibility intervals; the constrained
LASSO adds $\mathrm{P}_2$ and $\mathrm{D}_2$. Chapter~\ref{ch:homotopy} completes
the table.

\section{The homotopy}

Continuity of $x(\lambda)$ across breakpoints---the free block is uniquely
determined on each segment, and adjacent segments agree at the breakpoint---lets a
\emph{homotopy} recover the whole path: from a vertex with known active set, follow
\eqref{eq:affine} to the first event, update the active set by that one event,
repeat.
\index{homotopy}

\begin{algorithmbox}[The parametric active-set homotopy]
\label{alg:homotopy}
\begin{enumerate}
\item Compute a starting vertex and its active set (Section~\ref{sec:first} below).
\item While $\lambda$ has not reached its terminal value:
  \begin{enumerate}
  \item solve \eqref{eq:border} for two right-hand sides against one factorisation,
        obtaining $\alpha$, $\beta$, and the affine multipliers;
  \item compute every candidate step $t$ in the four families of
        Section~\ref{sec:events}, discarding inadmissible ones;
  \item take $t^\star = \min t$; if none is admissible, stop;
  \item apply the single event attaining $t^\star$ (Bland tie-break), update
        $(\Free,\Bound,\mathcal{S})$ and the factorisation, and record the
        breakpoint $x(\lambda_k + \sigma t^\star)$.
  \end{enumerate}
\end{enumerate}
\end{algorithmbox}

\subsection{The first turning point}
\label{sec:first}

The walk needs a starting vertex. For the portfolio instance, at $\lambda = +\infty$
variance is irrelevant and~\eqref{eq:pqp} degenerates into ``earn the most expected
return the bounds and budget allow''. With a single all-ones budget row that vertex
is computed greedily: initialise every weight at its lower bound, sort assets by
descending expected return, and walk down that order raising each weight toward its
upper bound until the budget is met. The last asset touched is only partially filled
and is the single \emph{free} asset of the first turning point.

Two details of that construction repay attention, because both are the kind of thing
a derivation omits and an implementation cannot.

First, the budget test needs a tolerance. The accumulated increment can only reach
$1$ up to rounding, and without slack the loop skips past the genuinely interior
asset and marks the next, zero-weight, asset as free.

Second, marking the last asset \emph{free} even when it lands exactly on a bound is
deliberate: it guarantees the first turning point has at least one free asset, so the
reduced system is nonsingular from the outset. Otherwise every asset sits on a bound
and the reduced solve degenerates.

For a general equality system the greedy fill no longer applies, and the
maximum-return vertex is instead the solution of a linear program,
$\max_x d\T x$ subject to $Ax = b$, $\ell \le x \le u$, which a simplex code
returns as a vertex. Only the first turning point needs the linear program; the loop
that follows is unchanged.

\subsection{Termination}

\begin{proposition}[Finite termination of the homotopy]
\label{prop:homotopy-finite}
Under the standing assumption, with a lowest-index tie-break, the homotopy visits
each active-set state at most once and terminates in finitely many steps.
\end{proposition}

\begin{proof}
There are finitely many active-set states, at most $2^{n+p}$. The parameter
$\lambda$ is monotone along the walk. The only failure mode is revisiting a state at
the same $t^\star = 0$, and Bland's lowest-index rule excludes that exactly as it
does for the simplex method.
\end{proof}

The bound $2^{n+p}$ is a finiteness statement, not an efficiency estimate. In
practice the number of breakpoints is empirically $O(n)$, and
Chapter~\ref{ch:reading} discusses what is and is not known about the gap.

\section{Two readings of the parameter}

Before Chapter~\ref{ch:homotopy} identifies the instances, one conceptual point.
The parameter $\lambda$ in~\eqref{eq:pqp} multiplies the \emph{linear} term while
the quadratic term carries no coefficient of its own. In the portfolio reading it is
therefore a \emph{return tilt}, not a risk aversion; sweeping it from $0$ (minimum
variance) to $\infty$ (maximum return) traces the whole efficient frontier. In the
regression reading it is a penalty, and the sweep runs the other way, from
$\lambda_{\max}$ (everything zero) to $0$ (the unpenalised fit).

Either way it plays the role a statistician reserves for a regularisation
parameter: a single knob weighing two competing goals, whose whole trajectory---not
any one point of it---carries the information. That the trajectory happens to be
computable exactly, rather than sampled, is what makes the subject worth a part of
this book.

\begin{remark}[Vector parameters]
\label{rem:vector}
Nothing forces $\lambda$ to be a scalar. If the parameter is a \emph{vector}---the
state of a plant, in the control instance---then the interval of breakpoints becomes
a polyhedral partition of parameter space into \emph{critical regions}, and the
homotopy becomes region traversal. That is multi-parametric quadratic programming
\cite{bemporad2002,tondel2003}. It is a genuine generalisation of the scalar device
rather than another instance of it, and this book treats only its scalar slice;
Chapter~\ref{ch:homotopy} says where the boundary falls.
\end{remark}

\begin{notes}
The piecewise-linearity of the solution path under a quadratic loss and a polyhedral
penalty was characterised within statistics by Rosset and Zhu \cite{rosset2007}; the
same fact is the piecewise-affinity of the mpQP solution over critical regions in
control \cite{bemporad2002}, and Roll \cite{roll2008} traces single-parameter paths
in that setting, explicitly generalising the LARS construction. The bordered-solve
presentation here follows \cite{schmelzer2026perspective}, and the first-turning-point
construction and its two details follow the \texttt{cvxcla} software paper
\cite{cvxcla}; the degenerate case in which every starting asset sits on a bound is
not addressed by the otherwise thorough step-by-step review of Markowitz and
Todd \cite{markowitz2021}.
\end{notes}

\begin{exercises}

\begin{exercise}
Verify that the coefficient matrix of~\eqref{eq:border} is nonsingular under the
standing assumption, adapting the proof of Proposition~\ref{prop:saddlenonsing}.
Where exactly is $H_{\Free\Free} \succ 0$ used, and where is full row rank of
$C_\Free$ used?
\end{exercise}

\begin{exercise}
Show that $\beta_\Bound = 0$ on every segment, and interpret this: what does it say
about a blocked variable as $\lambda$ moves?
\end{exercise}

\begin{exercise}
Take $n = 2$, $H = I$, $d = (1,0)\T$, bounds $0 \le x \le 1$, and the single equality
$x_1 + x_2 = 1$. Trace the path from $\lambda = \infty$ to $\lambda = 0$ by hand,
identifying every breakpoint and which of the four event families produced it.
\end{exercise}

\begin{exercise}
Show that for a genuine box bound the $\mathrm{D}_1$ event reduces to $\rho_i = 0$,
and that for the $\ell_1$ penalty it is $|\rho_i| = \lambda$. Why does the second
case make the active set ``induced'' rather than ``imposed''?
\end{exercise}

\begin{exercise}
Prove that between consecutive breakpoints the objective value
$\tfrac12 x(\lambda)\T H x(\lambda)$ is a quadratic function of $\lambda$, and that
the parametric objective $\tfrac12 x\T Hx - \lambda d\T x$ is concave in $\lambda$
along the path.
\end{exercise}

\begin{exercise}
Explain why the first turning point of the greedy construction has at least one free
asset, and construct data for which the naive version (no tolerance on the budget
test) marks the wrong asset free.
\end{exercise}

\begin{exercise}
Show that the number of breakpoints is at least the number of distinct active sets
visited, and give an instance with $n = 3$ where an asset leaves the free set and
later re-enters.
\end{exercise}

\begin{exercise}[Implementation]
Implement Algorithm~\ref{alg:homotopy} for the box-and-budget case (events
$\mathrm{P}_1$ and $\mathrm{D}_1$ only). Validate it against a general-purpose QP
solver called at $50$ values of $\lambda$: the traced segment, evaluated at each
value, should agree to solver precision. Report the largest discrepancy.
\end{exercise}

\begin{exercise}[Implementation]
Extend your tracer to inequality rows ($\mathrm{P}_2$, $\mathrm{D}_2$) and validate a
group-exposure cap the same way. How many extra breakpoints does one cap introduce?
\end{exercise}

\end{exercises}
